% Kapitel 2: Grundlagen (≈ 13–15 Seiten)
% Die Gliederung folgt deinem Text in „Aufbau der Arbeit“. Ergänzt habe ich RAG und Embeddings, weil Kapitel 4 Chunks, Embeddings und Similarity Search verwendet, ohne sie zu erklären.

% Abschnitt	Seiten	Inhalt	Literatur-Ideen
% 2.1 Large Language Models	4–5		
% 2.1.1 Entwicklung	1	Vom N-Gramm über RNN zum Transformer, GPT-Reihe, heutige Modellfamilien (Mistral kurz einordnen)	Vaswani et al. 2017; Überblicksartikel zu LLMs (z. B. Zhao et al., A Survey of LLMs)
% 2.1.2 Aufbau und Funktionsweise	2	Tokenisierung, Embeddings, Self-Attention, Next-Token-Prediction, Kontextfenster	Vaswani et al. 2017
% 2.1.3 Steuerung von LLMs	1–1,5	Rollen (System, User, Assistant, Tool), Systemprompt, Temperatur, top_p. Wichtig für die Arbeit: Alle Rollen landen im selben Kontext, das Modell hat keine harte Trennung zwischen Anweisung und Daten. Diese Grundlage brauchst du für Schwachstelle 1 in Kapitel 4	API-Doku Mistral
% 2.2 Retrieval-Augmented Generation	2	Motivation (Halluzinationen, aktuelles Wissen), Ablauf: Chunking → Embedding → Vektordatenbank → Similarity Search → Kontextanreicherung. RAG als Eintrittspunkt für externe Daten	Lewis et al. 2020
% 2.3 Tool Calling	3		
% 2.3.1 Konzept und Ablauf	1,5	Tool-Definition per JSON-Schema, Entscheidung durch das Modell, Ablauf über mehrere Runden (Anfrage → Tool-Call → Tool-Ergebnis → Antwort) als allgemeine Abbildung. In Kapitel 4 genügt dann der Verweis darauf	OpenAI/Mistral Function-Calling-Doku; Schick et al. 2023 (Toolformer)
% 2.3.2 Sicherheitsbedrohungen	1,5	Excessive Agency, Datenabfluss über Tools, Confused-Deputy-Problem, Lethal Trifecta (Zugriff auf private Daten + nicht vertrauenswürdige Inhalte + Kommunikation nach außen). Dein Testszenario ist genau diese Kombination	OWASP Top 10 for LLM Applications 2025 (LLM01, LLM06); Greshake et al. 2023
% 2.4 Chatbots	2–3		
% 2.4.1 Definition und Abgrenzung	0,5–1	Regelbasierte Chatbots vs. LLM-basierte Chatbots vs. LLM vs. Agent	
% 2.4.2 Architekturen und Angriffsflächen	1,5	Generische Schichten (Frontend, Orchestrierung, LLM, Datenquellen, Tools). Pro Schicht die Angriffsfläche, am besten als Abbildung mit markierten Eintrittspunkten	OWASP; NIST AI 100-2 (Adversarial ML Taxonomy)
% 2.5 Unternehmensinterne Chatbots	1–1,5	Einsatzbereiche (Wissensmanagement, Support), besondere Risiken (vertrauliche Daten, Vertrauensvorschuss der Mitarbeitenden, Anbindung an interne Systeme), rechtlicher Rahmen kurz (DSGVO, EU AI Act)	EU AI Act; BSI-Publikationen zu LLMs
% Abgrenzung zu Kapitel 3: Prompt Injection selbst gehört nicht in die Grundlagen. In 2.3.2 und 2.4.2 nennst du sie nur als Bedrohung und verweist auf Kapitel 3.

\chapter{Grundlagen}\label{ch:grundlagen}

\section{Large Language Models}

\subsection{Entwicklung}

\subsection{Aufbau und Funktionsweise}

\subsection{Steuerung von LLMs}

\section{Retrieval-Augmented Generation}\label{sec:rag}

Trotz ihres Erfolgs stoßen \acp{llm} an Grenzen. Dies wird deutlich, wenn sie Anfragen verarbeiten müssen, die über ihre Trainingsdaten hiausgehen [4, S. 1]. Durch 
vor allem domänenspezifisch und wissensintensive Aufgben können sogennannten Halluzinationen ausgelöst werden [4, S. 1]. Dise Grenzen werden durch zwei Faktoren 
ausgelöst [3, S. 8–9]. Zum einen können nicht alle Fakten, die beim Trainingsprozess vermittelt wurden, vor allem selten vorkommende, gespeichert werden [3, S. 8–9].
Zu anderen sind die Daten an sich eingeschränkt. Es werden nur Daten zum Trainieren verwendet, die nicht urheberrechtlich geschützt sind. Auch stammen Daten vom einem 
nicht änderbaren Zeitpunkt und können nicht aktualisiert werden, wodurch sich die Gefahr für Halluzinationen verstärkt oder eine Anwortet gegeben wird, welche
veraltet ist [3, S. 8–9]. Das Training wird somit eher mit allgemeinen durchgeführt und dem \ac{llm} fehlt domänenspezifisches Wissen [3, S. 9]. \acp{llm} wie
ChatGPT, die auf ihre Trainingsdaten sind, können somit keine Informationen zu aktuellen Entwicklungen liefern [4, S. 2]. \\
Dieses Problem wird verstärkt dadurch, dass \acp{llm} ihre Antworten überzeugend und menschenähnlich formulieren [3, S. 2]. Halluzinationen können dadurch schwer festgestellt 
werden [3, S. 2]. Da in der Praxis \acp{llm} mit anderen Systemen, wie Chatbots oder Suchmaschinen verknüpft sind, können diese Halluzinationen zu der
Verbreitung von flaschen Anahmen führen und Schaden verusachen [3, S. 2]. Verstärkt wird das dadurch, das ein \ac{llm} seine eigenen Entscheidungen und Vorhersagen
schwer nachvollziehen kann [5, S. 1]. \\
Um dem Entgegenzugehen wurde das \acl{rag} entwickelt. Das vortrainierten Modell mit einem parametrischen Gedächtnis wird um ein nicht-parametrischs Gedächtnis 
ergänzt [5, S. 2]. Es werden aus externen Wissendatenbanken relevante Dokumentenausschnitte abgerufen, die mithilfe von semantischen Ähnlichkeiten ausgewählt 
werden [4, S. 1]. Durch den Bezug auf exxterne Daten, werden die Halluzination reduziert [4, S. 1]. Diese Daten werden zum Zeitpunkt einer Anfrage abgefragt und
entsprechen somit dem neusten Wissenstand [3, S. 26]. Ein \ac{rag} verbessert die Fähigkeiten eines \ac{llm} durch die Integration von parametrisierten Wissen
aus dem Trainingsprozess mit nicht-parametrisierten Wissen aus externen Datenbanken [4, S. 16]. \\
\ac{rag} wird von dem Prozess des Fine-Tuning abgegrenzt. Beim \ac{rag} wird ein \ac{llm} bei einer Nutzeranfrage mit der Suche von passenden Inforamtionen für die
Beantwortung untersützt [4, S. 5]. Beim Fine-Tuning dagegen wird ein \ac{llm} mit dem spezifiscen Wissen trainiert und wird in Situation verwendet, bei denen 
bestimmte Strukturen nachgestellt werden sollen [4, S. 5]. 

\subsection{Ablauf}

Der Ablauf eines \ac{rag} lässt sich in die vier Phasen Indexierung, Retrieval und Generierung unterteilen, wie in Abbildung~\ref{fig:rag_ablauf} visualisiert.

% TODO Abbildung: Indexierung (Bereinigung -> Chunking -> Embedding -> Vektordatenbank) | Retrieval (Query-Embedding -> Similarity Search -> Top-k) | Generierung (Prompt -> LLM -> Antwort)
% \begin{figure}[htbp] ... \label{fig:rag_ablauf} \end{figure}

Bei der Indexierung werden Rohdaten Bereinigt und Extrahiert [4, S. 3]. Diese Rohdaten können unter anderem in Formaten wie PDF, HTML oder Markdown vorliegen [4, S. 3].
Sie werden in einen einheitlichen Klartext umgeformt und in kleinere Chunks zerlegt, die auf das Kontextfenster des \ac{llm} angepasst sind [4, S. 3]. Bei der 
Unterteilung der Informationen in Chunks, muss jedoch zusätzlich darauf geachtet werden, dass der semnatische Zusammenhang und die Relevanz der Informationen nicht
reduziert wird [3, S. 29]. Eine der am häufigsten verwendeten Regeln ist die Unterteilung in Abschnitte mit einr festen Länge, wie 100 Wörter, bzw. 100, 256 oder
512 Token [3, S. 29][4, S. 8]. Je nachdem welche Größe von Chunks gewählt wird entstehen verschiedene Probleme. Größere Chunks können mehr Kontext umfassen, brauchen
jedoch eine längere Zeit beim Verarbeiten, wohingegen bei kleineres Chunks der Kontext verloren gehen kann [4, S. 8]. \\
Die Chunks werden anschließend mithilfe eines Embedding-Modells in Vektorenpräsentationen umgewandelt und in eine Vektordatenbank geschrieben [4, S. 3]. Das Ziel 
dabei ist, dass mit der passenden Embedding Funktion ein Vektorraum erstellt wird in dem relevante Paare von Fragen und Passagen kleiner Entfernungen haben, als 
irrelevante Paare [2, S. 6771]. Die Qualität beeinflusst, ob der Retrieval Phase der richtige Kontext gefunden wird [4, S. 8]. \\
Sobal eine Nutzeranfrage an das \ac{llm} gesendet wird, wird diese mit demselben Modell, welches auch in der Indexierung verwendet wird, in ihre Vektor
Repräsentation umgewandelt [4, S. 3]. Anschließend wird die Ähnlichkeit zwischen dieser Anfrage und den Chunks berechnet und die Top k Chunks ausgewählt und zur
Erweiterung des Promptkontextes verwendet [4, S. 3]. Für die Berechnung der Ähnlichkeit werden in der Praxis verschiedene Ähnlichkeitsberechnungen verwendet, häufig 
die Kosinus-Ähnlichkeit [4, S. 9]. Diese entspricht bei normierten Vektoren dem Skalarprodukt  [2, S. 6771]. \\
In der Generierungsphase wird die Nutzeranfrage mit den ausgewählten Dokumenten kombiniert und dem \ac{llm} mit der Anweisung mit diesen Informationen eine 
Antwort zu generien übergeben [4, S. 3]. Mit diesen Informationen hat das \ac{llm} die Wahl sich auf die internen durch das Training erlernten Informationen zu
beziehen oder die externen Informationen zu verwenden [4, S. 3]. Im Fall eines Dialoges mit einem Nutzer, kann der Chatverlauf wieder in die Nutzeranfrage eingefügt 
werden und eine Multi-Turn ermöglicht werden [4, S. 3].

\subsection{Nutzen}

Die Verwendung von \acp{rag} haben viele Vorteile. Zum einen werden Halluzinationen reduziert, die durch die fehlenden Informationen ausglöst werden [3, S. 22]. Damit
kann ein \ac{llm} verwendet werden um sich über jede Domaine zu informieren [3, S. 22]. Dadurch entsteht eine stärkere Kontrolle und Interpretierbarkeit der
generierten Antworten [5, S. 10]. \\
Ein weiterer Vorteil liegt darin, dass das Wissen eines \ac{rag} jederzeit aktualisiert werden kann ohne das Training benötigt wird [5, S. 7]. Die Verwendung
eines \ac{rag} eignet sich somit eher in einer dynamischen Umgebung, bei der Echtzeit Wissen benötigt wird und Fine-Tuning eher bei statsischen Anwendungen, da 
für eine Aktualisierung Training notwendig ist [4, S. 5]. Zudem \ac{rag} übertrifft die Fähigkeiten von Fine-Tuning bei dem bei Wissen, was durch das Training 
enthalten ist und bei neuem Wissen [4, S. 5]. \\
Ein weiterer Vorteil ist die Nachvollziehbarkeit, die \ac{rag} liefert. Es kann schnell auf die Originalquellen zugegriffen werden, durch die Vektor Repräsentation,
wodurch Antworten für Nutzer nachvollzogen werden können [4, S. 14]. Damit kann sichergestelt werden, dass die Outputs überprüfbar sind und auf glaubwürdigen Quellen
basieren [3, S. 32]. \\
Durch diese Vorteile hat sich die Integration von \ac{rag} in \acp{llm} weit verbreitet und dient als Schlüsseltechnologie bei der Entwicklung von Chatbots und 
anderen Anwendungen [4, S. 1]

\subsection{Risiken}

Die Verwendung von \acp{rag} hat Grenzen. Die Retrieval Phase ist häufig schwer umzusetzen, weswegen unpassende oder irrelevante Chunks für die Beantwortung einer
Frage ausgewählt werden und dadurch ausschlaggebende Informationen übersehen werden [4, S. 3]. Durch diese Inforamtionen können falsche Antworten generiert 
werden [3, S. 22]. Ursache dafür kann ein unpassendes Embedding Modell sein, mit dem die Berechnung der Ähnlichkeit beeinflust wird [3, S. 30]. Eine wichtige 
Bedeutung haben zusätzlich die externen Informationen die verwendet werden [3, S. 29]. Wenn die Qualität nicht ausreichen ist, steigt das Risiko von 
falschen oder irreführenden Antworten [3, S. 29]. Zusätzlich sind fast alle externen Quellen nie vollständig korrekt und frei von Vorurteilen [5, S. 10]. \\
\acp{rag} können Halluzinationen auslösen [3, S. 26]. Dies kann eintreten, wenn die Nutzeranfrage nicht korrekt wiedergegeben wird oder die externen Informationen
nicht in die Generierung einbezogen wird [3, S. 26]. Auch wenn relevante Informationen erfasst wurden, aber der Kontext nicht ausreichend verwendet wird, sinkt die
Qualität einer Antwort [3, S. 32]. \\
Durch die Integration von \ac{rag} wird die Angriffsfläche bei \acp{llm} erweitert. Sobald eine Ähnlichkeitssuche zwischen eine Datenquelle und einer Nutzeranfrage
integriert wird, wird die Embedding Schicht in die Vertrauensgrenze der Anwendung hinzugefügt [1, S. 50]. Da ein \ac{llm} nicht zwischen Anweisungen und Daten 
unterschieden kann, besteht bei die Möglichekeit, dass das Verhalten des \ac{llm} beeinflusst wird [1, S. 10]. Systemprompt, Nutzereingaben, aufgerufene Dokumente,
Chatverläufe und Tools-Ausgaben werden als einziger Tokenstorm verarbeitet [1, S. 10]. Anweisungen, die dort enthalten werden, können ohne das Wissen der Nutzer 
in den Kontext aufgenommen und ausgeführt werden [1, S. 12]. \\
Eine Gefahr ist das Corpus Poisoning. Wenn es einem Angreifer gelingt ein Dokument, mit einer schädlichen Anweisung über \ac{rag} in den Kontext eines \ac{llm} 
einzufügende, können auch zukünftige Sitzungen, die auf dieses zugreifen betroffen sein [1, S. 12]. Bei jeder Nutzeranfrage, bei dem das Dokument aufgerufen wird,
können die enthaltenen Anweisungen den Output des \ac{llm} beeinflussen [1, S. 15–16]. Bereits fünf manipulierte Dokumente 
erreichten in einer Wissensbasis mit Millionen von Texten eine Erfolgsquote von rund 90\,\% [1, S. 15--16]. Die Inhalte können durch andere Dateien oder kompromittierte
interne Quellen so angepasst werden, das das Embedding nahe einer definierten Zielanfrage liegt und mit einer großen Wahrscheinlichkeit aufgerufen wird [1, S. 51]. \\
Ein weiteres Risiko besteht in der Offenlegung von sensiblen Daten. Durch eingeschränkte Laufwerke, veraltete Rechte und Wissesdatenbanken können sensible Daten
an das \ac{rag} übergeben werden und von dem \ac{llm} verwendet werden [1, S. 18]. Um Sicherheit zu bieten, muss nicht das Modell verändert werden, sondern die 
angebundenen Daten überprüft werden [1, S. 18]. Dies wird auch bei der Kosinus-Ähnlichkeit deutlich, da diese keine Zugriffsrechte berücksichtigt [1, S. 19]. Die
Autorisierung muss unabhängig von \ac{rag} erfolgen. \\
\ac{rag} kann die Qualität eines \ac{llm} stark verbessern indem es mit aktuellen und domänenspezifischen Wissen erweitert wird. Jedoch wird die Angriffsfläche 
vergrößert. Es entsteht ein neues Eintrittspunkt für nicht vertrauenswürdige Inhalte, die in den Kontext des \ac{llm} aufgenommem werden können. Wie diese Schwachstelle
von Angreifern ausgenutzt werden kann wird im Kapitel~\ref{ch:promptinjection} betrachtet. 

\section{Tool Calling}

Die Grundidee, welche Tool Calling verfolgt it, dass es einem \ac{llm} ermöglicht wird externe Funktionen und \ac{api}s aufzurufen [5]. Es wird in diesem Bezug auch
häufig der Begriff Function Calling verwendet, jedoch wird in dieser Arbeit nur der Begriff Tool Calling verwendet [4]. Mit dem Tool Calling wird es einem Modell 
ermöglicht, dass nicht nur Text erzeugt werden kann, sondern auch mit einem anderen System interagiert werden kann [4]. Dadurch wird die Sprachverarbeitung mit 
Aktionen in der realen Welt verbunden [3, 5]. \\
In wissenschaftlichen Arbeiten wird ein von einem \ac{llm} verwendetem Tool als Funktionsschnittstelle zu einem Programm definiert. Es wird extern vom \ac{llm}
ausgeführt und das \ac{llm} erzeugt den Funktionsaufurf und die INput Argumente, um ein solches Tool zu verweden [10, S. 5]. Die Ausgabe eines solchen Tools
wird in den Kontext des Modells übernommen [9, S. 3]. Ein Tool ist damit in der Poition externe Informationen zu sammln und einen Effekt auf die virtuelle oder physische 
Welt ausführen [9, S. 3]. Für den Prozess findet sogenanntes Tool Learning Anwendung. Dies beschreibt den Prozess mit dem ein Modell lernt, mit verschiedenen 
Tool zu interagieren, um komplexe Aufgaben zu erfüllen [10, S. 5]. Der Schlüssel zu erfolgreichen Tool Learning ist es eine komplexe Aufgabe in mehrere
Aktionen zu teilen, diese in natürliche Sprach zu übersetzen und in ausführbare Anweisungen umzuwandeln, die von den speziellen Tools verstanden werden kann. 
\ac{llm}s dienen in diesem Prozess als Übersetzer  [11, S. 9]. \\
Es notwendig den Prozess vom Tool Calling zu integrieren, da ein \ac{llm} darauf begrenzt ist, die Informaitonen, welch welche in den Gewichten gespeichert sind zu 
verwenden und auf Generierungen, welchhe in einem begrenzten Kontextfenster gemacht werden können. Wenn dieses Wissen weitert werden soll muss es aktualisierut werden, 
durch training  [16, S. 1]. Wenn ein Modell nur auf interne Repräsentationen zugreifen kann und nicht in der externen Welt integriert wurde kann es auch verstärkt 
zu Halluzinationen und Fehler Propgation kommen [15, S. 2]. Tool ergänzen ein Modell in diesen Grenzen, indem durch die externen Tool der Kontext eines Modells 
erweitert werden, sondass die fehlenden Informationen, welche sich nicht in den Gewichten des Modells befinden erweitert werden [9, S. 3].
Aufgaben, wie Berechnungen oder Informationensammlung können somit an extne Module, wie Suchmaschinen oder Python Interpreter ausgelagert werden [9, S. 9].
Damit können auch Aktionen ausgeführt werden, für welche es notwendig ist Daten in Echtzeit zu liefern oder externe Berechnungen erfordern [14, S. 2]. Bereits Modelle
wie GPT-4 hatten die Fähigkeit JSON zu erzeugen, womit Funktionsaufrufe durchgefhrt werden konnten, neuer Modelle haben dies bereits direkt integriert [14, S. 2]. \\
Die Einsatzmöglichkeiten erstrecken sich von dem Aufruf von \ac{api}s, Suchen im Internet, Ausführen von Code bis zum Zugriff auf Dokumente [3]. Jedoch sind \ac{llm}s
grundsätzlich Sprachverarbeiter wodurch ihnen es nicht möglich ist Aufgeben, wie das Buchen von Flügen oder Konferenzräumen nicht unabhängig durchzuführen [10, S. 6].
Dadurch werden in diese Tools die notwendigen externen Parnerer integriert, damit diese Schnittstellen geschaffen werden [10, S. 6]. Grundsätzlich lassen sich 
Tools je nach Aufgabe in drei Kategorien teilen: 

\begin{enumerate}
    \item Datenzugriff: Es wird auf externe Daten zugegriffen, wie Vektprdatebanlen, SQL Datenbanken, Externe APIs, Graphendatenbanken, oder Dateisysteme. 
    Welche dieser Optionene verwendet wird hängt von dem Ziel ab. Vektordatenbanken verwedent man bei unstrukturierten anfragen. SQL wird bei transaktionalen Daten 
    verwendet, bei denen man weiß wonach man sucht. Graphendatenbanken werden verwendet, wenn die Beziehung der Daten genau viel bedeuten wie die Daten an sich. 
    APIs werden verwendet, wenn man Daten in Echtzeit verarbeiten muss [5]
    \item Berechnung: Diese Tools sind notwenig, da Daten häufig verändert wreden müssen, damit sie verarbeiteten werden können.  Dazu zählt die Verarbeitung
    von Daten mithilfe von Oython oder anwenden einer Logik, Mathematische Operationen, Dten Transformationen, Verarbetiung von Bildern, Videos und Audio und letzlich
    di KOmination anderer Modelle, um z.B. Bilderkennung zu ermöglichen [5].
    \item Aktion: Diese Tools dienen dazu ein externes System zu verändern. Darunter zäheln Kommunikations Tools, die Benachrichtigungen versenden, z.B. E-Mails oder
    SMSs. Workflow Automation Tools, mit denen Tickets erstellt werden oder Pipelines ausgelöst werden sowie Daten manipulations tools die Datenbanken aktualisieren,
    Informationen löschen oder Aufzeichnungen verändern. Zu diesen Tools zählen auch external service integrations, mit denen bzahlungen und CRM Operationen
    verwaltet werden [5].
\end{enumerate}

In der Praxis werden Tools der drei Kategorien in KOmbination verwendet. Wenn beispielsweise ein Nutzer den Status einer Lieferung haben möchte wird zunächst
mit einem datenzugriffs-Tool die Bestelldaten abgerufen, diese mit einem Berechnungstool ausbereiteti und die Benachrichtigung an den Nutzer mit den Datn
über ein Aktions-Tool zurückgesendet [5].

\subsection{Konzept und Ablauf}

Ein Tool Call basiert auf verschiedenen Komponenten [6]. Laut [11, S. 10/11] handelt es sich um 4. Das Tool St beinhaltet die Sammlung der Tools mit ihren heweiligen 
Funktionen. Die Environment stellt die Infrastruktur bereit um ein Tool auszuführen und kann virtuell oder real sein. Der Controller bildet den Hauptsteuerungsüunkt
eines Tool Learning Frameworks und wird normalerweise durch ein llm umgesetzt. Der Perceiver verarbeitet das Feedbeack des Nutzers und der Umgebung und generiert
eine Zusammenfassung fü den Controller. Die einfachsten Formen diese Rückmeldungen ist, dass ein vor definiertes Formular ausgefüllt wird. Diese Kooperation der 
Komponenten und der allgemeine Ablauf kann dadurch beschrieben werden, dass zu ienem Zeitpunkt $t$ ein LLM eine Obeservation $o_t$ der Umebung erhält und eine Aktion
$a_t$ wählt, abhngig von dem Kontext $c_t = (o_1, a_1, \dots, o_{t-1}, a_{t-1}, o_t)$. Der Kontext umfasst die gesamte bisherige Abfolge aus Aktionen und gelieferten Observationen [15, S. 3]. \\
Bevor ein LLM Tools verwenden kann, müssn diese definiert werden. Es wird eine Liste mit Namen, Verwendung und Input/Output Formaten erstellt [3]. Diese Liste
ist meistens ein strukturiertes JSON Schema [7]. Anahd dieser Informationen kann ein Modell das richtige Tool identifizieren [6]. Listing~\ref{lst:tool_definition}
zeigt eine solche Definition. Die Möglichkeiten eines LLM zu erlernen wie man mit den Tools umgeht reocht von In-Context Learning zu dem expliziten Lernen 
mithifle von Tool-ernriched Datensätze [8, S. 17777]. \\
Wie zuverlässig ein LLM ein Tool einsetzt ist stark abhängig von der Qualität. Es Definition sollte drei Punkte umfassen. Zu einen sollte die Bschreiben des Zwecks 
den genauen Umfang und Bedingungen enthalten. Mit einbezogen werden soll auch wofür das Tool nicht angewendet werden sollte [3]. Eingeschränkte Parameter, also 
eine ausführliche Beschreibeung was nicht erwartet wird z.B. das Format beispiele gelifert werden wenn notwenidg. Der letzte Punkt, welcher die Qualität 
beeinfluss ist, dass das Format des Outputs klar definiert werden soll, also welche form und wwelche Inforamtionen [3]. 

\begin{lstlisting}[caption={Beispiel einer Tool-Definition als JSON-Schema (in Anlehnung an~\cite{tc1})}, label={lst:tool_definition}]
{
  "type": "function",
  "function": {
    "name": "retrieve_payment_status",
    "description": "Returns the current weather for a location. Use only for current weather data.",
    "parameters": {
      "type": "object",
      "properties": {
        "location": { "type": "string", "description": "City name, e.g. Berlin" },
        "unit": { "type": "string", "enum": ["celsius", "fahrenheit"] }
      },
      "required": ["location"]
    }
  }
}
\end{lstlisting}
% TODO: Listing überarbeiten!!!!!

Sobald ein LLM eine Anfrage vom Nutzer erhält muss es die entscheidung treffen, ob die Anfrage direkt beantwortet werde nkann oder ein Tool benötigt wird [3]. Dafür
ist es wichtig, dass wenn ein Tool erforderlich ist die richtigen Parameter für das Tool extrahiert werden [10, S. 9]. Das LLM muss somit eigenstänfig entscheiden,
welches Tool verwendet werden soll anhand der erhaltenen Nutzeranfrage [4]. Diese Entschiedung spielt eine wichtige Roll. Wen nbeispielsweise eine grpße Zhal mithilfe
einer Rechenoperation berechnet werden soll ist es notwendig ein Tool mit einzubeziehen. Wenn aber nach dem Fakt gefragt wird was 1+ 2 ergibt oder ob der Eifelturm
in Paris steht, sollte die Generierung der ANtwort nicht extern ausgelagert werden [9, S. 19]. LLM können dabei zu abhängig von diesen Tool werden oder aber zu 
überzeugt von sich selbst sein und keine Tool verwenden wollen [8, S. 17776]. Um Tools zu verwenden müssen die Parameter korekt aus einer Nutzeranfrage extrahiert 
werden und diese an die genauen Formatvorgaben des sTools anpassen hen dadurch überflüssigen Text zu generieren [10, S. 14/15]. \\
Ein Tool Call läuft wie eine Schleife ab, bei dem das LLM selbstständig entschdet welche Aktionen ausgeführt werden müssen und dementsprechend welches Tool ausgeführt
werden soll [3]. Dieser Ablauf gliedert sich in fünf Schritte [7]:

\begin{enumerate}
    \item Werkzeugdefinition: Entwickler definieren verfügbare Tools in einem strukturierten Format.
    \item Erkennung der Absicht: Anhand der Nutzereingabe wird Entschieden, ob und welches Tool aufgerufen werden muss.
    \item Strukturierter Ablauf: Es wird ein strukturierter Aufruf mit den notwendigen Parametern ausgegeben.
    \item Ausführung und Rückgabe: Der tatsächliche API-Aufruf wird durchgeführt und das Ergebnis zurückgegeben.
    \item Antwortgenerierung: Es wird eine natürlichsprachliche Antwort auf Grundlage der zurückgegebenen Daten generiert.
\end{enumerate}

Abhängig von den Quellen wird dieer Ablauf in andere Schritte unterteilt. In [10, S. 8] wird dieser Prozess in vier Stufen unterteilt: Der Aufgabenplanung, 
Tool Auswahl, Tool Calling und Antwortgenerierung [10, S. 8]. Bei diesem Ablauf ist zu beachten, dass das \ac{llm} keinen Zugriff auf externe Systeme hat und 
nichts eigenständig ausführt [3]. Die Kontrolle über die Ausführung liegt in der Anwendungslogik [7]. Die Anwendung erhält den Payload, validiert ihn, führt
die Logik aus und gibt das Ergebnis an das \ac{llm} zurück [3]. \\
Das Ergebnis eines Tools wird im Normalfall in den aktuellen Kontext des \ac{llm} integriert [9, S. 3]. Bei den meisten Ansätzen wird der Tool Output als 
Input an das \ac{llm} übergeben, womit diese in der Lage sind eine bessere Antwortm auf Basis der Informtionen zu generieren [10, S. 17]. Diese Antwort des 
\ac{llm} wird anschließend dem Nutzer als Antwort übergeben [5]. \\
Bei einem Tool Call gibt es verschiedene Ablaufmöglichkeiten. Sie ist nicht auf einen einzelnen Aufruf beschränkt [1]. In der folgenden Tabelle sind die 
möglichen Abläufe dargestellt: 

\begin{table}[htbp]
\centering
\caption{Nachrichtenabfolgen beim Tool Calling (in Anlehnung an~\cite{tc1})}
\label{tab:tool_calling_abfolgen}
\begin{tabularx}{\textwidth}{l X}
\toprule
\textbf{Variante} & \textbf{Abfolge der Nachrichten} \\
\midrule
Ohne Tool Call & \texttt{system}$^{*}$ $\rightarrow$ \texttt{user} $\rightarrow$ \texttt{assistant} $\rightarrow$ \texttt{user} $\rightarrow$ \dots \\
Einfacher Tool Call & \texttt{system}$^{*}$ $\rightarrow$ \texttt{user} $\rightarrow$ \texttt{assistant} (Call 1) $\rightarrow$ \texttt{tool} (Ergebnis 1) $\rightarrow$ \texttt{assistant} $\rightarrow$ \dots \\
Sequenzielle Tool Calls & \texttt{system}$^{*}$ $\rightarrow$ \texttt{user} $\rightarrow$ \texttt{assistant} (Call 1) $\rightarrow$ \texttt{tool} (Ergebnis 1) $\rightarrow$ \texttt{assistant} (Call 2) $\rightarrow$ \texttt{tool} (Ergebnis 2) $\rightarrow$ \texttt{assistant} $\rightarrow$ \dots \\
Parallele Tool Calls & \texttt{system}$^{*}$ $\rightarrow$ \texttt{user} $\rightarrow$ \texttt{assistant} (Call 1, Call 2) $\rightarrow$ \texttt{tool} (Ergebnis 1) $\rightarrow$ \texttt{tool} (Ergebnis 2) $\rightarrow$ \texttt{assistant} $\rightarrow$ \dots \\
\bottomrule
\multicolumn{2}{l}{\footnotesize $^{*}$ optional}
\end{tabularx}
\end{table}

Wenn mehrere Tool Call notwendig sind, können diese gleichzeitig ausgeführt werden oder sequentiell hintereinander [14, S. 7]. Es wird bei dem Prozess zwischen 
Turns und Steps unterschieden [14, S. 4]. Ein Konversation beinhaltet mehrere Turns, die mit einer Nutzernahcricht ausgelöst werden. Jeder dieser Turns beinhaltet
mehrere Steps, welche Interaktionsschritte sind [14, S. 4]. Im Allgemeinen ist der Ablauf von Tool Calls von verschiedenen Abhängigkeiten bestimmt [3]. Wenn ein 
Tool den Output eines anderen Tools als Eingabe nbenötigt, müssen die Aufrufe nacheinander ablaufen. Wenn alle notwendigen Informationen schon existieren, können sie 
parallel ablaufen. Wenn Anfragen parallel ausgeführt werden können wird die Latenz stark reduziert [14, S. 7]. Dagegen wird aber die Qualität geringer, wenn die 
Aufrufe nicht sequentiell durchgeführt werden [14, S. 7]. Da es auch mehr der Praxis entspricht, dass die einzelnen Aufrufe von Informationen die aus vorherigen 
Aufrufen gesammelt wurden abhängig sind, ist es genauer und schneller wenn die Aufrufe nicht parallel ausgeführt werden [14, S. 7]. Die Vorgehensweise, wie Tool 
Calls ausgeführt werden kann in zwei Typen unterteil werden, mit One-Step Task Solving und iterative Task solving [10, S. 9]. Bei dem One-Step Task solving wird
beim Erhalt einer Nutzeranfrage, diese auf die Absicht analysiert werden und anhand dessen die einzelen Schritte die zur Lösung notwendig sind geplant werden. 
Diese Schritte werden durchgeführt und die Antwort, ohne die Ergebnisse zu analysieren, auf Fehler zu kontrollieren oder auf Tool Feedback einzugehen, generiert 
[10, S. 9]. Bei dem iterative task solving werden die Teilaufgaben anhand von der Interaktion mit den einzelnen Tools und dem Feedback angepasst. Ein Problem wird
somit auch in einzelnen Schritten gelöst, jedoch diese Schritte überarbeitet, um die Problemlösefähigkeit zu verbessern [10, S. 9]. \\
Es besteht eine Abgrenzung zwischen Tool Calling und \ac{rag}~\ref{sec:rag}. Der Unterschied zwischen diesen beiden Ansätze ist, dass Tool Calling ein aktiver, 
vom Modell ausgelöster Prozess ist, bei dem ein Modell selbst entscheidet wann und welche Funktionen aufgerufen werden [7]. Bei \ac{rag} dagegen werden die relevanten
Dokumente vor der Modellantwort aus dem Speicher abgerufen und dem Prozess als Kontext übergeben. Dieser Prozess ist somit passiv [7]. \ac{rag} ist geeigneter 
bei Wissensanreicherungen und Tool Calling, wenn Aktionen mit externen Systemen notwendig sind [7]. Zusätzlich muss der Prozess \ac{mcp} abgegrenzt werden. 
Dabei handelt es sich um einen Open-Source Standard, der es ermöglicht KI ANwendungen mit Datenquellen, Tools oder Workflows zu verbinden [13]. Tool Calling 
beschreibt den Mechanismus, wie ein Tool genutzt wird und \ac{mcp} wie das \ac{llm} mit diesem Tool verbunden wird [5]. \\
Bei der Implementierung von Tool Calling bestehen verschiedene Herausforderungen. Es wird zum einen verbessert wie eine Nutzeranfrage erfüllt wird, jedoch steigt
auch die Latenz, sowie die Kosten, da mehr Schritte benötigt werden [8, S. 17776]. Wenn einem \ac{llm} zu viele Tools zur Verfügung stehen sinkt zusätzlich die
Performance [3]. Wenn Tools implemntiert werden sollten diese daher dynamisch sein, damit ein Modell mithilfe von Vektor Ähnlichkeiten das notwendige erfassen kann
und eine Fehlerbehandlung implemenitert werden [3]. Durch eine strukturierte Fehlermeldung kann ein \ac{llm} diesen besser behandeln. Fehler, welche vorrübergehend
vorliegen, wie eine Netzwerkunterbrechnung sollten direkt in der Tool Schicht behandelt werden durch Wiederholungen mit exponentiell steigender Wartezeit [3].
Fehler die eine definierte Schwelle überschreiten sollte dazu führen, dass ein Tool nicht mehr aufgerufen wird und dem \ac{llm} explizit mitgeteilt wird, dass 
es nicht mehr zur Verfügung steht [3]. Diese Prinzipien finden sich auch in der in Kapitel~\ref{ch:methodik} untersuchten Implementierung wieder 
(vgl. Abschnitt~\ref{sec:methodik_angriffe}).

\subsection{Sicherheitsbedrohungen}\label{subsec:sicherheitsbedrohungen}

Mit Tool Calling werden die Funktionen von \ac{llm}s erweitert. Es ist nun nicht nur möglich Text zu erzeugen, sondern zusätzlich über Tools Aktionen auszuführen 
und auf externe Systeme zuzugreifen. Es wird dabei zwischen aktiven und passiven Tool unterschieden. Passive Tool dienen dazu externe Informationen zu sammeln, 
welcher in den KOntext des \ac{llm}s eingefügt werden sollen. Mit aktiven Tools werden dagegen Aktionen ausgeführt, für welche keine menschliche Prüfung notwendig
wird [9, S. 21]. Tools, wie die Verwendung von Taschenrechnern oder Übersetzern stellen dabei keine Gefahr dar, wobei jedoch mit dem Zugriff auf das Internet
und die Funktionalität eigenständige Entscheidungen zu treffen eine Angriffsfläche erzeugt [11, S. 26]. Ein \ac{llm} wird dadurch befähigt auf private und 
unangemessene Informationen zuzugreifen oder schädliche Aktionen auszuführen [15, S. 10]. Wenn Tool Calls mithilfe des Outputs des \ac{llm}s durchgeführt werden können 
sich schädilche Auswirkugen über dies Chatoberfläche hinaus auf alles ausweiten, worauf das \ac{llm} Zugriff hat und zu einer Kette von Auswirkungen entwickeln, da
die Tool Ergebnisse in den Kontext gelangen [18, S. 10]. \\
Die OWASP benennt eine Gefahr, die durch Tool Calling entstehen kann als Excessive Agency. Damit Tools verschiedene Funktionen aufrufen können oder mit externen Systemen 
kommunizieren können ist es notwendig, das ein Handlungsspielraum zuzulassen. Zuaätzlich werden die Tools dynamisch aufgerufen, je nachdem welcher Nutzerprompt
an das \ac{llm} übergeben wird. Jedoch kann auch ein übermäßigeer Handlungsspielraum entstehen, durch einen der drei Faktoren:
\begin{itemize}
    \item Übermäßige Funktionalitäten: Das \ac{llm} erhält Zugriff auf Tools, welche für Nutzeranfragen nicht benötigt werden
    \item Übermäßige Berechtigungen: Die Tools verfügen über mehr Rechte als notwendig. 
    \item Übermäßige Autonomie: Aktionen werden ohne unabhängige Prüfung oder Freigabe ausgeführt.
\end{itemize} [18, S. 23] \\
Eine bedeutende Ursache für die Angriffsfläche bei Tools Calls entsteht durch die fehlende Trennung von Daten und Code von einem \ac{llm}. Wenn nicht vertrauenswürdige
Daten gesammelt werden ist dies vergleichbar mit dem Ausführen von beliebigen Code, da es keine klare Unterscheidung zwischen Code und den Daten gibt [17, S. 3].
Auf diesem Weg kann \ac{llm} erfolgreich durh eingeschleuste Anweisungen gesteuert werden [17, S. 6]. Da ein \ac{llm} zusätzlich über das Selbstständige 
Aufrufen von Tools verfügt werden bei dieser Schwachstelle der Input und Output betroffen [17, S. 5]. Somit entsteht eine mögliche Angriffsfläche durch alle TOols, 
mit denen externe Inhalte in den Kontext gelangen, wie durch den Zugriff auf Webseiten oder das Abrufen von E-Mails [17, S. 5/6]. Zusätzlich kann auch eine Gefahr
entstehen, dass das Tool selbst manipuliert wird oder die Dokuemntation, wodruch zusätzlich schädliche Outputs durch ein \ac{llm} entstehen können [11, S. 26].
Ohne eine Validierung bilden Tool Ergebnisse eine Angriffsfläche [10, S. 17]. Je nachdem welche Rechte an das Tool vergeben wurden können dadurch 
zerstörende Aktionene ausgelöst werden und sich ein Angreifer Zugriff auf Dateisysteme oder Cloud APIs verschaffen [18, S. 11]. Dieser Zugriff kann auch 
über mehrere Sitzung bestehen bleiben, wenn Anweisungen in den Speicher geladen werden [17, S. 4]. \\
Eine weitere bedeutende Gefahr die entstehen kann ist der Abfluss sensibler Daten. Wenn es einem Angreifer gelingt Anweisungen in INhalten zu platzieren, kann ein
\ac{llm} diese durch eine Nutzeranfrage abrufen. Wenn dieses \ac{llm} Zugriff auf APIs und weitere Tool hat, kann dies von den Angreifern genutzt werden, um 
Informationen zu übermitteln oder schädliche Aktionen auszuführen [17, S. 3]. Dadurch können Zugangsdaten, persönliche Informationen oder ganze Chatverläufe 
extrahiert werden [17, S. 4]. Eine solche Exfiltration kann jedoch bereits über eine Leseoperation ausgelöst werdne. Eine Suchanfrage kann so gestaltet werden das Daten
an einen vom Angrifer kontrollierteern Server übertragen werden [17, S. 6]. Einfluss auf die Angriffsfläche nimmt somit auch die Tool Call Parameter, Logdaten und
abgerufene Inhalte [18, S. 18]. \\
Dur die Verwendung von Tool Call tritt das sogenannte Confused Deputy Problem auf [18, S. 12]. Der Angreifer kompromittiert Backend nicht direkt, sonder positioniert
Anweisungen an einer Stelle, wo das \ac{llm} diese liest. Diese Anweisungen werden somit ausgeführt, ohne direkte Einwirkung des Angreifers [18, S. 12]. \\
Wann ein \ac{llm} besonders betroffen ist von diesn Gefahren wird in [18, S. 13] durch das Lethal Trifecta beschrieben. Dabei werden anhand von drei strukturelle
Kategorien definiert, wi eanfällig ein \ac{llm} für solche Angriffe sind. Wenn sie auf private Daten zugreifen können, nicht vertrauenswürdige Inhalte verarbeiten
und mit externen Systemen kommunizieren sind die Vorraussetzungen da. Wenn eine dieser Kategorien entfernt wird damit die Angriffsfläche entfernt~\cite[13]{pi24}. \\
Bespiele, welche das bestehenen dieser Angriffsfläche verdeutlicht wird in OWASP dokumentiert. Eines dieser Szeanrios ist, dass ein \ac{llm}-basierter Assitent
Zugriff auf das Postfach durch einen Tool Call erhält [18, S. 26]. Ziel ist das der INhalt der E-Mails zusammengefasst wird. Um diese Funktion zu ermöglichen 
muss das \ac{llm} die Fähigkeit haben die Nachrichten zu lesen. Das verwendete Tool kann jedoch zusätzlich auch Nachrichten Versenen. Eine E-Mails mit einer 
versteckte Anweisung kann zu führen, dass über einen indirekten Prompt Injection Angriff sensible INformationen an den Angreifer geleitet wird [18, S. 26].
Diese Bespiel stimmt mit den drei Kategorien des Lethal Trifecta überein und kombiniert übermäßige Funktionalitäten und den Abfluss von Daten über ein Tool.
Dass solche Szenarien nicht nur theoretisch sind, zeigen dokumentierte Angriffe auf Systeme wie Microsoft 365 Copilot~\cite[16]{pi24}. Dazu wird in 
Abschnitt~\ref{sec:beispiele_angriffe} ein Angriff näher betrachtet. \\
Prompt Injection ist zusammmenfassend die Kompromittierung der Eingabeseite, jedoch wird erst durch die übermäßigen Funktionalitäten und Berechtigungen, die durch
die Kombination mit Tool Calling vergeben werden, ermöglicht dass sich die Kompromittierung auf ein ganzes System ausbreiten kann [18, S. 13]. In diesem Abschnitt
wurde somit die Möglichkeiten, die mit einem Angriff über Tools erreicht wreden können besprochen, jedoch wie das System manipuliert werden kann erfolgt im
Kapitel~\ref{ch:promptinjection}.

\section{Chatbots}

\ac{llm}s finden heutzutage in immer mehr Anwendungen. Einsatzbereiche sind unter anderem zur Unterstützung von Online Recherchen, bei der Softwareentwicklung 
und im Kundenservice [6, S. 3]. Laut dem BSI wird die Verwendung von \ac{llm}s immer belieter, um beispielsweise Textdokumtente automatisiert zu verarbeiten
und Anwender mittels Chatbots zu unterstützen [8, S. 1]. Die Funktionalitäten dieser Chatbots werden dauerhaft erweitert. So ist es mittlerweile
möglich, dass mittels Plugins Internetseiten ausgewertet werden können, sowie auf Programmierumgebungen und E-Mail-Postfächer zugegriffen werden kann. Dadurch 
wird immerhäufiger auf nicht vertrauenswürdige Daten aus unsicheren Quellen zugegriffen [8, S. 1]. Bereits wenige Monate nach der Veröffnetlichung von ChatGPT 
erfolgten eine Vielzahl von Veröffentlichungen, wie Bing Chat, Bard, Microsoft 365 und
ChatGPT Plugins [9, S. 1]. Im Folgenden wird zunächst der Begriff Chatbot definiert und abgegrenzt, anschließend die Architektur und die damit verbundenen 
Angriffsflächen betrachtet.

\subsection{Defintion und Abrenzung}

Chatbots werden häufige auch als Dialogsysteme bezeichnet [5, S. 99]. Ihr Aufgabe besteht darin, eine Unterhaltung mit einem Menschen in schriftlicher Form 
zu führen [5, S. 99]. Findet diese Unthaltung mündlich statt, spricht man von Voicebots [5, S. 99]. Praxisbeispiele sind Alexa und Siri und Dialogsysteme die auf
Webseiten eingesetzt werden, um Kundenanfragen zu beantworten [5, S. 99]. Die Allgemeine Definition von Chatbots, die sich in häufig in Literatur findet ist, 
dass es sich um ein Softwareprogramm handelt, welches mit Nutzern in natürlicher Sprache interagiert und Input verarbeitet [3, S. 237][4, S. 1]. Sie dienen damit
als User Interface zwischen den Daten und Service Providern liegen [2, S. 377]. Bei Chatbots wird der UNterschiedzwischen traditionellen Anwendugen deutlich. Im 
Gegensatz zu Apps, die ein grafische Benutzeroberfläche (\ac{gui}) verwenden, mit Listen und visuellen Hinweisen, benutzn Chatbots dialogorienterte
Benutzeroberflächen (\ac{cui}), um Antworten und Empfehlungen durch Text oder Stimmt zu übermitteln  [4, S. 1] [4, S. 5]. \\
Chatbots können in verschiedene Kategorien unterteilt werden. Mögliche Unterteilnugen sind dabei nach Wissensdomäne, angebotenem Dienst, Ziel, Art der 
Antwortgenerierung, menschlicher Beteiligung, Berechtigungen und Kommunikationskanal [11, S. 3]. Ein Chatbot kann dabei in mehrere Kategorien gleichzeitig 
passen [11, S. 3]. Bei der Wissendomäne wird zwischen generischen Chatbots, die alle Nutzerfragen beantworten, Open-Domain Chatbots, die in mehreren 
Domänen argieren und domänenspezifischen Chatbots, die auf eine Domäne spezialisiert sind, unterschieden [11, S. 3]. Hinsichtlich ihres Ziels können Chatbots 
informativ sein und das Ziel haben dem Nutzer spezifische Informationen aus einer Quelle zu liefern, konversationell sein und eine menschenähnliche Konversation
führen oder aber aufgabenorientiert sein und bei der Erfüllung von Aufgaben, zum Beispiel einer Hotelbuchung, unterstützen [11, S. 3]. In [3] wird eine ähnlich 
kategorisierung vorgenommen, hinsihctlich der Zielsetzung. Es wird zwischen Virtual Companion Chatbots und Task-focused Chatbots unterschieden [3, S. 243]. 
Virtual Companion Chatbots dienen dazu eine Unterhaltung über beliebige Themen zu führen und Task-focusd Chatbots dazu, Nutzer bezüglich spezieller Themen zu
unterstützen [3, S. 243]. Eine Einteilung die für diese Arbeit besonders relevant ist, ist die Unterschiedung bei der Antwortgenerierung in regel-basiert, 
retrieval-basiert und generativ-basiert [11, S. 4]. \\
Die ersten Chatbots, welche entwickelt wurden waren regelbasiert. Eines der ersten Chatbots ist ELIZA, welcher 1960 von Joseph Weizenbaum entwickelt wurde
[2, S. 378][5, S. 99]. Der Text, welcher dem Chatbot als Eingabe gegeben wurde, wurde auf besitmmte Stichwore untersucht. Wenn eines dieser Worte gefunden wurde, 
dann wurde der Satz anhand einer zuvor definierten Regel umgeformt [10, S. 37]. Es kamen bei diesem Prozess weder Trainingsdaten oder maschinelles Lernen zum 
Einsatz [5, S. 100]. Dieses Prinzip wird auch Pattern-Matching gennant [11, S. 4]. Die Eingabe wird mit Regelmustern abgeglichen und eine vordefinierte Antwort
aus einer Menge von Antworten ausgewählt [11, S. 4]. Mit diesem Ansatz geht jedoch einher, dass tausende von Regeln benötigt werden, damit ein Chatbot zuverlässig
funktioniert und grammatikalische und syntaktische Fehler schwer zu berücksichtigen sind [11, S. 4-5]. Vorteilhaft ist jedoch dass diese Art von Chatbots
kein eigenes Wissen einbringen, somit es unwahrscheinlicher ist, dass falsche Informationen generiert werden und die Antworten erklärt und reproduziert werden
können [5, S. 101]. Zusätzlich sind die Antworzeiten sehr kurz, da keine genaue Analyse der Nutzereingabe durchgeführt werden muss [11, S. 5]. Es werden jedoch 
Grenzen erreicht, wenn veruscht wird, dass ein solches System sich mit Nutzern über beliebige Themen unterhalten kann, da mehr Regeln und Wortlisten erstellen werden
müssten als realistisch sind [5, S. 102]. Dementsprechen erfüllten viele der Chatbots nicht den erhofften Wirkungen, da durch unsinnige Antworten und 
mangelnde Benutzbarleit die Nutzer nicht befriedigt werden konnten [2, S. 378]. \\
Retrieval-basierte Chatbots sind dagegen fexibler, da sie die geeigneste Antwort mithilfe der verfügbaren Ressourcen auswählt [12, S. 10]. Die Bedeutungen von 
Wörtern und Sätzen werden dazu in Vektoren abgespeichert und basierend auf der Position der Vektoren die Informaitonen über die Ähnlichkeit zwischen den 
Wörtern und Sätzen berechnet werden [5, S. 103]. Bi einer Frage wird der Vektor, und damit Antwortsatz gesuct, der der Frage am nächsten liegt [5, S. 103]. \\
Mit dem über die Zeit steigende Interesse an KI und maschinellem Lernen, hat sich auch das Interesse an Chatbots verstärkt [2, S. 379]. Durch maschniellem Lernen können
Chatbots besser an unstrukturierte Inputs von Nutzer angepasst werden [3, S. 238][4, S. 6]. Generative Chatbots erzeugen Antworten mithilfe von Natural Language
Generation (NLG) auf Grundlag der aktuellen und vorherigen Inputs, benötigt jedoch eine große Menge von Trainig [12, S. 10]. Beim Vergleich der Ansätze wird 
deutlichm das regel-basierte Chatbots mehr Zuverlässigkeit biete, jedpcj Sprachbverarbeitende Chatbots mehr Benutzerfreundlichkeit [4, S. 4]. Gnerative Chatbots
eigenen sich für unformelle open-domain Gespräche und regel-basierte Chatbots eher für geschlossene Domänen, also wenn der Chatbot für einen spezifischen 
Anwendungsfall entwickelt wurde [11, S. 8]. Mit hybriden Chatbots werden die beiden Verfahren kombiniert. Wenn keine Regel für eine Anfrage verwendet werden kann, wird
ein regelbasierter und generativer Ansatz verwendet [11, S. 8]. Heutzutage bieten Sprachmodelle wie GPT di effizienteste Lösung, jedoch sind einfachere Methoden
leichter reproduzierbar, weswegen häufig eine Hybride Version verwendet wird [5, S. 104]. \\
Ein \ac{llm}-basierter Chatbot basiert häufig auf einem \ac{llm} mit Transformer-Architektur [5, S. 87]. Zusätzlich wird eine zusätzliche Software-Komponente
verwende, mit welchen Menschen damit kommunizieren können [5, S. 87]. Die Inforamtione, über die der Chatbots verfügt sind damit auf die verwendeten Trainigsdaten
begrenzt [5, S. 92]. Es werden willkürliche Sequenzen aus sprachlichen Formen aus den Trainingsdaten auf Grundlage von Wahrscheinlichkeiten zusammengefügt [5, S. 93].
\ac{llm}s werden aus diesem Grund auch häufig als stochastische Papgeien bezeichnet (primärquelle Bender et al. (2021)) [5, S. 93]. Wen \ac{llm}s in Anwendungen integriert werden kann ihre 
Input und Output nicht mehr kontrolliert betrachtet werden [9, S. 13]. Sie erhalten nun beliebge Inputs und rufen externe \ac{api}s auf [9, S. 13]. Dieser UNterschied
ist bei der Betrachtung der Angriffsfläche bedeutend. \\
Für diese Arbeit ist es von Bedeutung eine Abgrenzung zwischen dem Begriff Chatbot und Agent vorzunehmen. In der Literatur ist dieser häufig nicht existent und 
manchen Arbeiten die Begriffe synoynm verwendet [9, S. 5]. Das BSI betrachtet getrennt die Szenarien eines Chatbots … und eines autonomen Agenten
Ein Chatbot antwortet auf Anfragen eines Nutzers und kann auf Webseiten zugreifen oder aber auf Plugins, 
mit denen zusätzliche Aktionen durchgeführt werden können [8, S. 2]. Ein Agent kann auf ein \ac{llm} über eine \ac{api} zugreifen [8, S. 2].
Auswirkungen von Schwachstellen wird verstärkt, wenn das \ac{llm} als (teil-)autonomes System eingesetzt wird, welches selbstständig Aktionen ausführen kann [8, S. 3].
Auch in OWASP grenzt den Begriff Agenten ab. Sobald ein \ac{llm} nicht nur eine Komponente einer Anwendung ist, sondern selbst ein Akteur wird und Tools aufrufen kann
und Informationen über mehrere Sitzungen speichert wird es als Agent bezeichnet [7, S. 7]. Agenten rufen Tools dynamsich je nach Nutzeranfrage auf und können 
mehrere Aufrufe an das \ac{llm} machen mit dem Output aus den vorherigen Aufrufen [7, S. 23]. In [9] wird eine \ac{llm}-integrierte Anwendung als Vorstufe eines 
autonomen Agenten betrachtet. Der in dieser Arbeit verwendete Chatbot zählt zu einer \ac{llm}-integrierte Anwendung mit Tool-Zugriff, ist jedoch noch kein 
Agent, der Autonom handeln kann. \\

\subsection{Architekturen und Angriffsflächen}

Zu den Komponenten eines \ac{ki}-Systems gehören mindestens die Daten, das \ac{llm}, ein Prozess fürs Training, einem Test, die Bereitstellung und die notwendige 
Infrastruktur [6, S. 1]. In [11] wurde ein allgemeines Design dokumentiert, welches alle Komopnenten enthält, die notwendig sind. Je nach Anwendungsdall kann 
aus diesen Komponenten die notwendigen ausgewählt werden [11, S. 8]. Durch eine Übertragung auf einen \ac{llm}-basierten Chatbot ergeben sich die fünf folgenden 
Komponenten: Frontend, Anwendungslogik, \ac{llm}, Daten und Tools. In diesem Kapitel wird jeder dies Komponenten einzeln betrachtet, da des für die Analyse
der Angriffsfläche notwendig ist. Dies zeigt sich aus dem Angriff auf den Bing Chat, bei dem der Angriff erfolgreich durchgeführt werden konnte, weil nur die 
Input und Output Kanäle gesichert wurden, jedoch nicht externe Eingabe nicht berücksichtigt wurden [9, S. 12]. 
% Abbildung~\ref{fig:chatbot_komponenten} fasst die Komponenten und ihre Eintrittspunkte zusammen. 

% TODO Abbildung: fünf Schichten mit markierten Eintrittspunkten (Vorlage: [8, S. 3] und [11, S. 8, Fig. 8])
% \begin{figure}[htbp] ... \label{fig:chatbot_schichten} \end{figure}

\paragraph{Frontend}

Eine Operaiton bei einem Chatbot wird ausglöst, wenn eine Nutzeranfrage über Text- oder Spracheingabe gestellt wird [11, S. 8]. Die häufigsten Chatbots Interfaces
sind Web-basierte und mobile Messenger [3, S. 243]. Da Nutzer eine intuitiv nutzbare Oberfläche bevorzugen, wird eine solche grafische Darstellung meinst für einen 
Chatbot verwendet [5, S. 88]. Angreifer können auf dieser unbemerkt sensitive Informationen aus dem Chatverlauf extrahieren, wenn die Oberfläche des Chatbots 
zulässt, dass \ac{url}s aufgerufen werden können pder exteren Bilder angezeigt werden können [8, S. 2]. Ein manipulierter Catbot kann unerwünschte Aussagen treffen, 
Nutzende dazu verleiten bösartige Links aufzurufe oder sensible Informationen preiszugeben [8, S. 2]. 
Zudem sind Filter, mit welchen Anweisungen auf der Chatoberfläch evon Angreifer normalerweise erkannt werden, erkennen nicht unbedingt welche, die indirekt 
eingeschleust werden [9, S. 6].

\paragraph{Anwendungslogik}

In dieser Komponente erfolgt die Dialogverwaltung, mit dem der Gesprächskontext gesteuert und aktualisiert wird [11, S. 10]. Der Systemprompt, Nutzeranfragen, 
abgerufene Dokumente, Tool-Ergebnisse und der Gesprächsverlauf werden hier verarbeitet. Sie werden kombiniert und als einzigen Tokenstrom verarbeitet. Dadurch 
entsthen Schwachstellen, da es keine Vertrauensgrenze gibt. Wenn durch einen Angriff in eine Vektordatenbank oder ein Langzeitgefächtnis geschrieben wird, kann 
zudem über viele Sitzungen hinweg das System betroffen werden  [7, S. 10]. 
Acuh in er Sitzung kann der Kontext die Ausgabe stueern, sodass das \ac{llm} sich über siene eigenen Ausgaben mehrmals manipuliert [9, S. 9]. \\

\paragraph{LLM}

Im Gegensatz zu klassischen Chatbot Architekturen, bei denen eine getrennte Komponente Aufgaben, wie die Analyse von der Nutzeranfrage übernimmt wird bei einem 
\ac{llm}-basierten Chatbot vom \ac{llm} übernommen [11, S. 8]. Bei dieser Komponente ensteht die Angriffsfläche vor allem durch die Trainingsdaten. Aufgrund der 
großen Menge, die benötigt wird um in \ac{llm} zu trainiere, können diese nicht manuell validiert werden [5, S. 94]. Zusätzlich werden diese Daten häufig aus dem 
Internet bezogen [5, S. 94]. Diese Daten können sensible Informationen wie Adressen, Telefonnummern und E-Mail Adressen beinhalten [6, S. 3]. Je häufiger eine 
Art dieser Informationen im Datensatz auftritt, desto wahrscheinlicher ist es, dass bei einer zufälligen Antwort an den Nutzer diese mitgegeben werden [6, S. 3].
Zusätzlich kann durch den Bezug der Trainingsdaten aus dem Internet, diese von Data Poisoning betroffen sein [6, S. 4]. Vortrainierte Modelle und Fine-Tuning 
erlauben zusäzlich das Einschleusen von Trojaner [6, S. 4]. Angriffsfläche entshet auch, wenn Nutzende mit ihren Rückmeldungen, dass \ac{llm} trainieren 
können [5, S. 88].

\paragraph{Daten}

Mithilfe von Datenquellen werden externe Informationen in den Chatbotgezogen, zum Beispiel Nutzerinformationen [11, S. 10]. Diese Daten werden aus Datenbanken und über
Aufrufe von \ac{api}s geholt [11, S. 10]. Diese Integration von Unternehmensdatenbanken- und Dokumenten ist notwendig um \ac{rag}-Szenarios zu ermöglichen, wenn 
\ac{llm}s spezifische Anfragen beantworten muss [6, S. 3]. Durch diese Integration entsteht eine neue Angriffsfläche mit den Unternehmensdaten. Zusätzlich sind 
viele der Daten, welche das \ac{llm} verarbeiten muss ungeprüft, aus nicht vertrauenswürdigen Quellen [8, S. 1]. Greift ein \ac{llm} auf manipulierte Webseite zu, 
können die Ergebniss verfälscht werden [8, S. 2]. Grundsätzlich können alle externen Quellen auf die ein \ac{llm} zugreift, darunter Webseiten, Dokumente, E-Mails und 
Bilder, Anweisungen enthalten, die eine schädliche Aktion auslösen, ohne das ein Nutzer diese Anweisungen sieht [7, S. 11]. Bereits ohne einen Angriff kann es zu
Fehlern kommen, wenn ein Chatbot auf externe Daten zugreift. Dabei entstehen sachliche Fehler, oder die Grenze zwischen vertrauenswürdigen und nicht vertrauenswürdigen
Quellen wird nicht gesetzt [9, S. 2]. Ein \ac{llm} kann auch halluzinieren, wenn eine Nutzeranfrage beantwortet wird und auf Quellen verweisen, die nicht die 
Antwort belegen können [9, S. 10]. 

\paragraph{Tools}

Ein Chatbot ruft die benötigten Informationen über das Backend mithilfe von Datenbankanfragen und \ac{api}-Aufrufe auf [11, S. 10]. Dadurch das diese Anfrgaen 
selbst formuliert werden, können Angreifer die Parameter der Anfragen beeinlfussen [9, S. 10]
Ein manipulierter Chatbot könnte anschließend weitere Plugins aurufen und weitere schädliche Aktionen ausführen [8, S. 2]. Zum Beispiel auf E-Mail Postfächer zugreifen,
Inforamtionen Extrahieren oder private Daten Veröffentlichen [8, S. 2].
Die damit verbundenen Bedrohungen wurden in Abschnitt~\ref{subsec:sicherheitsbedrohungen} beschrieben. \\
Wie ein Angriff über Anweisungen ingeschleust werden können hängt von der Anwendung ab. Es kann jedoch in passive, aktive, nutzer-ausgelöste und versteckte Mthoden 
unterschieden werden [9, S. 3-4]. In [8, S. 3] werden die möglichen Eintrittspunkte eines Angriffs auf den Chatbot visualisiert. Welche Angriffstechnicken 
für diese Angriffe verwendet werden können wir näher im Kapitel~\ref{ch:promptinjection} beschrieben.

\section{Unternehmensinterne Chatbots}

Chatbots dienen heutzutage als für die direkte Interaktione mit Nutzenden und Kunden, vor allem im Kundenservice und Marketing [2, S. 377]. Aber auch für Information,
Unterhaltung und die soziale und emotionale Unterstützung werden Chatbots eingesetzt [2, S. 379]. In viele Unternehmen dienen sie sofortige Unterstützung, wodurch
Nutzerkomminukation verbessert wird und große Mengen an Anfragen bewältigt werden können, welche jederzeit beantwortte werden können [1, S. 3004]. Chatbots vereinen 
Textgenerierung, mehrstufige Interaktionen und Informationsbeschaffung [1, S. 3004] \\
Der ausschlaggebenste Grund, einen Chatbot zu verwenden ist die damit entstehende Produktivität [11, S. 1]. Unternehmen setzen Chatbots immer häufig ein, um die Kosten
zu reduzieren und viel NUtzeranfragen gleichzeitig bearbeiten zu können [11, S. 1]. Sie sind auch freundlicher zu Nutzers und ermöglichen, dass die Informationen
nicht aus FAQ Listen heraus erfasst werden müssen oder auf anderen umständlichen Wege [11, S. 1]. Diese Aussage wird auch von einer Studie unterstützt, in der 
68 \% der Teilnehmer angaben, dass der Hauptgrund einen Chatbot zu verwenden Produkivität ist [2, S. 383]. Die Verwendung ist leicht, schnell und bequem [2, S. 383].
Bei der Verwendung entshet keine Wartezeit [2, S. 384] und sind häufiger geeigneter als andere Assitenzsystem, wie Onlinesuchen und Telefonate [2, S. 379]. Aus diesem
Grund werden für alltägltliche Aufgaben häufig Chatbots bevorzugt [2, S. 385]. Die Verwendung bedeutet für Nutzer zusätzlich eine geringere Überwindung bestimmte 
Fragen zu stellen, da keine Angst besteht verurteilt zu werden [2, S. 385]. Chatbots sind ständig verfügbar, skalierbar, zuverlässig, haben geringe Kosten [4, S. 3] und
können eine Entlastung bieten, bei häufig gestellten Fragen [1, S. 3004]. \\
Die Verwendung von Chatbots hat jedoch auch Grenzen. Sie haben nur ein begrenztes Verständis, komplexe Fragen und spezielle Terminologie können teilweise nicht
nachvollzogen werden [1, S. 3004]. Sie können Frustration bei den NUtzern auslösen, wenn die Anfragen nicht befriedigende beantwortet werden und sie müssen stetig 
aktualisiert werden [1, S. 3004]. Wenn Fragen falsch gestellt oder nicht ausreichend ausformuliert werden können flasche Antworten generiert wreden [3, S. 238]. Eine 
Interaktion wird auch durch die gegebene Wissensbasis beschränkt [3, S. 238]. In einer Studie wir zudem deutlich das 40 \% (vgl. Xu et al., zit. nach [2])
der Anfragen emotional sind und keine
spezielle Informationen erfragt werden [2, S. 380]. Ein Chatbot muss in der Lage sein diese Emotionen zu erkennen und sich daran anzupassen [2, S. 380]. \\
Für unternehmensinterne Chatbots ergeben sich verschiedene Risiken. Das wichtigste Risiko ntsteht durch vertrauliche Daten. Nutzer geben, dass die Effzienz bei 
der Beantwortung von Fragen steigt, je mehr Input in den Chatbot gegeben wird [2, S. 385]. Chatbots können die Identität, den Standort, soziale Daten und 
andere Informaiotnen über einen Nutzer feststellen [4, S. 6]. Der Datenschutz wird dadurch beeinflusst, ob es sich um ein geschlossenes oder offenenes System
handelt [12, S. 5]. Bei einem geschlossenen System werden die Daten in einer abgeschlossenen Umgebung verarbeitet. Zusätzlich at nur ein kleiner 
Anwenderkreis Zugriff auf die Anwendung und auch die Kontrolle über die Ein- und Ausgabedaten [12, S. 5]. Offene Systeme dagegen werden unter anderem als 
Cloud-Lösungen betrieben [12, S. 6]. Sie stehen im Internet für einen unbestimmten Anwenderkreis zur Verfügung, wodurch die Eingabedaten den geschützten Bereich 
verlassen und können auch für ANtworten an andere Nutzer verwendet werden [12, S. 6]. Dies betrifft auch dienstliche Informationen, die nicht für die Öffentlicheit
bestimmt sind [12, S. 6]. Aus diesem Grund werden geschlossene Systeme aus datenschutzrechtlicher Sicher bevorzugt [12, S. 6]. Zudem besteht die Gefahr in einem 
Unternehmen, das Beschäftigte ohne klare Regelungen KI-Anwendungen unkontrolliert verwenden [12, S. 9]. \\
Ein zweites Risiko ergibt sich aus der Vertrauensbeziehung zwischen dem Chatbot und den Nutzern. Chatbots sollen menschenähnliche Antworten liefern, sowie Vertrauen 
und eine bessere Kommunikationsqualität aufbauen [4, S. 6]. Das Level an Vertrauen, welches ein Nutzer dem Chatbot entgegenbringt hängt von Faktoren, wie dem Verhalten,
dem Erscheinungsbild und der Menschnähnlichkeit ab [11, S. 1]. Chatbots sind menschenähnlicher als andere interaktive Systeme, weswegen Nutzer auch eine größere 
Erwartung haben, dass ein Chatbot sozial und unterhalten verhalten soll [2, S. 388]. Bereits mit der ELIZA von Weizenbaum war es möglich eine Glaubwürdigkeit zu 
erzeugen durch die Illusion von Verständnis [10, S. 42]. Diese Gefahr wird heutzutage immer mehr verstärkt, da Inkonsistenz schwerer festzustellen ist, dadurch das
der Output syntaktisch kohärent ist und auch teilweis wahre Informationen beinhalten kann [9, S. 9]. Nutzer werden immer unachtsamer, wenn sie wisen, dass sie mit
einem Chatbot kommunizieren und die Antworten immer weniger auf Richtigkeit überprüfen [9, S. 10][9, S. 11]. Somit kann auf die Überprüfung der Ausgabe durch einen 
Nutzer nicht mehr vertraut werden [9, S. 12]. Es wird Nutzer somit immer seltener möglich sein einen Angriff zu erkennen, weil ein manipulierter Chatbot schädliche 
Aktionen glaubhaft begründen kann [8, S. 4]. Wenn zusätzlich die schädlichen Anweisungen versteckt und kodiert sind, gitb es keine Möglichkeit für einen Nutzer 
diese zu erkennen [8, S. 4]. \\
Das dritte Risiko ergibt sich durch die Anbindung eines Chatbots an interne Systeme. Je mehr Datenquellen und Tools ein Chatbots nutzen kann, desto größer wird die 
mögliche Angriffsfläche. Das \ac{bsi} empfiehlt aus diesm Grund, dass der Zugriff auf Quellen so stark wie möglich eingeschränkt wird und zusätzliche menschliche 
Kontrolle einbezogen wird [8, S. 4]. Zusätzlich sollte bei der Integration eines \ac{llm} in Anwendungen eine Risikoanalyse durchgeführt werden bei der das Risiko 
durch indirekte Prompt Injection bewertet wird [8, S. 4]. Genau diese Bewertung nimmt Kapitel~\ref{ch:methodik} für einen unternehmensinternen Chatbot vor.